% BhashaLLM Capstone Poster (Landscape)
% Compile with: pdflatex capstone_poster_landscape.tex
\documentclass[final]{beamer}
\usepackage[orientation=landscape,size=a1,scale=1.1]{beamerposter}
\usepackage{graphicx}
\usepackage{booktabs}
\usepackage{tikz}
\usepackage{amsmath}
\usepackage{array}
\usepackage{lmodern}
\usetikzlibrary{positioning, arrows.meta}

\definecolor{bhBlue}{RGB}{18, 64, 118}
\definecolor{bhTeal}{RGB}{19, 140, 126}
\definecolor{bhLight}{RGB}{235, 243, 250}
\definecolor{bhAccent}{RGB}{242, 111, 49}

\setbeamercolor{title}{fg=bhBlue}
\setbeamercolor{author}{fg=black}
\setbeamercolor{institute}{fg=black}
\setbeamercolor{block title}{fg=white,bg=bhBlue}
\setbeamercolor{block body}{fg=black,bg=bhLight}
\setbeamercolor{alertblock title}{fg=white,bg=bhAccent}
\setbeamercolor{alertblock body}{fg=black,bg=bhLight}
\setbeamercolor{headline}{bg=white}
\setbeamercolor{footline}{fg=white,bg=bhBlue}

\setlength{\columnsep}{2.2cm}
\setlength{\columnseprule}{0mm}

\title{\Huge \textbf{BhashaLLM:} বাংলা ভাষা ও হ্যান্ডরাইটিং বোঝার জন্য ক্যাপস্টোন সিস্টেম}
\author{\Large Student Name \quad \textbar \quad ID: 000000}
\institute{\Large Department of CSE \quad \textbar \quad University Name}

\begin{document}
\begin{frame}[t]
\begin{columns}[t]

% Column 1
\column{0.32\textwidth}
\begin{block}{সমস্যা ও প্রেরণা}
\vspace{0.2cm}
\begin{itemize}
  \item বাংলা একটি কম-রিসোর্স ভাষা; উচ্চমানের OCR ও LLM টুলের অভাব রয়েছে।
  \item হাতে লেখা বাংলা অক্ষরের ভ্যারিয়েশন বেশি হওয়ায় OCR নির্ভুলতা কমে যায়।
  \item শিক্ষার্থীদের জন্য স্বয়ংক্রিয় টিউটরিং ও ব্যাখ্যা-সহায়তা জরুরি।
\end{itemize}
\end{block}

\begin{block}{লক্ষ্য}
\vspace{0.2cm}
\begin{enumerate}
  \item হাতে লেখা বাংলা শব্দ সঠিকভাবে শনাক্ত করা।
  \item OCR আউটপুট থেকে শিক্ষামূলক ব্যাখ্যা/ফিডব্যাক তৈরি করা।
  \item অফলাইন ও লোকাল মডেল ভিত্তিক ব্যবহারযোগ্য API তৈরি করা।
\end{enumerate}
\end{block}

\begin{block}{ডেটাসেট ও ট্রেনিং}
\vspace{0.2cm}
\begin{itemize}
  \item OCR: বাংলা হস্তলিখন ডেটাসেট (গ্রাফিম-রুট, স্বরবর্ণ, ব্যঞ্জনবর্ণ)।
  \item LLM: বাংলা ইনস্ট্রাক্ট ফাইনটিউনিং ও কাস্টম অ্যাডাপ্টার।
  \item ট্রেনিং পাইপলাইন: ডেটা প্রিপারেশন $\rightarrow$ ট্রেনিং $\rightarrow$ ভ্যালিডেশন।
\end{itemize}
\end{block}

\begin{alertblock}{মূল অবদান}
\vspace{0.2cm}
\begin{itemize}
  \item লোকাল মডেল-ভিত্তিক বাংলা OCR + LLM স্ট্যাক।
  \item API-চালিত টিউটর ও দার্শনিক প্রতিক্রিয়া।
  \item অফলাইন ডিপ্লয়মেন্ট বান্ধব আর্কিটেকচার।
\end{itemize}
\end{alertblock}

% Column 2
\column{0.36\textwidth}
\begin{block}{সিস্টেম আর্কিটেকচার}
\vspace{0.4cm}
\centering
\begin{tikzpicture}[
  node distance=1.0cm,
  box/.style={rounded corners, draw=bhBlue, fill=white, thick, align=center, minimum height=1.1cm, minimum width=5cm},
  arrow/.style={-Latex, thick, color=bhBlue}
]
  \node[box] (input) {ইনপুট ইমেজ};
  \node[box, below=of input] (pre) {প্রি-প্রসেসিং \\ (কনট্রাস্ট, রিসাইজ)};
  \node[box, below=of pre] (ocr) {OCR মডেল \\ (ResNet-34 + Grapheme ক্লাসিফায়ার)};
  \node[box, below=of ocr] (text) {রিকগনাইজড টেক্সট};
  \node[box, below=of text] (llm) {LLM টিউটর/ফিডব্যাক};
  \node[box, below=of llm] (api) {FastAPI প্রতিক্রিয়া};

  \draw[arrow] (input) -- (pre);
  \draw[arrow] (pre) -- (ocr);
  \draw[arrow] (ocr) -- (text);
  \draw[arrow] (text) -- (llm);
  \draw[arrow] (llm) -- (api);
\end{tikzpicture}
\vspace{0.4cm}
\end{block}

\begin{block}{ইন্টারফেস ও API}
\vspace{0.2cm}
\begin{itemize}
  \item \textbf{/api/analyze} — OCR ইনফারেন্স + মেট্রিকস
  \item \textbf{/api/chat} — টিউটরিং ও ব্যাখ্যা
  \item \textbf{/api/philosophical} — দার্শনিক প্রতিফলন
\end{itemize}
\vspace{0.2cm}
\fbox{\parbox{0.9\linewidth}{\centering
ডেমো স্ক্রিনশট/ইন্টারফেস প্রিভিউ এখানে যোগ করুন}}
\end{block}

\begin{block}{মূল ফলাফল (উদাহরণ)}
\vspace{0.2cm}
\centering
\begin{tabular}{lcc}
\toprule
মেট্রিক & OCR & LLM \\
\midrule
Accuracy & 92.4\% & --- \\
CER      & 7.6\%  & --- \\
Response Quality & --- & 4.6/5 \\
\bottomrule
\end{tabular}
\vspace{0.2cm}
\end{block}

% Column 3
\column{0.32\textwidth}
\begin{block}{ইউজ-কেস}
\vspace{0.2cm}
\begin{itemize}
  \item বাংলা হ্যান্ডরাইটিং যাচাই ও শেখা
  \item শিক্ষার্থীদের জন্য স্বয়ংক্রিয় টিউটরিং
  \item বাংলা সাহিত্যিক কনটেন্ট বিশ্লেষণ
\end{itemize}
\end{block}

\begin{block}{ফিউচার ওয়ার্ক}
\vspace{0.2cm}
\begin{itemize}
  \item লাইন/প্যারাগ্রাফ লেভেল OCR
  \item মাল্টিমোডাল ব্যাখ্যা ও অডিও ফিডব্যাক
  \item মডেল কম্প্রেশন ও মোবাইল ডিপ্লয়মেন্ট
\end{itemize}
\end{block}

\begin{block}{টিম ও কৃতজ্ঞতা}
\vspace{0.2cm}
\begin{itemize}
  \item Supervisor: \textbf{Dr. Name}
  \item Team Members: [Name 1], [Name 2], [Name 3]
  \item বিশেষ ধন্যবাদ: Lab/Department
\end{itemize}
\end{block}

\begin{block}{যোগাযোগ}
\vspace{0.2cm}
\begin{itemize}
  \item \textbf{Email:} student.name@university-domain
  \item \textbf{GitHub:} github.com/Pronaaf2k/BhashaLLM
\end{itemize}
\vspace{0.2cm}
\fbox{\parbox{0.85\linewidth}{\centering
QR কোড এখানে দিন}}
\end{block}

\end{columns}
\end{frame}
\end{document}
